\subsection{Which evidence matters? Zones and lesion classes}
\label{sec:res-evablate}
The design of the evidence bottleneck rests on two assumptions that the preceding results do not test: that resolving lesions by retinal zone helps, and that all four lesion classes contribute. Table~\ref{tab:evablate} answers both questions by training the full \mname\ head with parts of the evidence tensor set to zero (including in the rule scores), using three seeds.

\paragraph{Zones add little.} Restricting the evidence to the whole-field zone gives a QWK of 0.736, compared with 0.745 for the reference with all six zones (a difference of $-0.009$), and removing the whole-field zone but keeping the four quadrants and the central disc gives 0.733 ($-0.012$). With seed-level standard deviations of 0.007--0.022, neither difference can be distinguished from noise. On DDR, then, \emph{the global lesion statistics carry practically all of the evidence that the zone-wise statistics carry}; the zone resolution that motivated the token design (Section~\ref{sec:rationale}) is not shown to be beneficial. We see two explanations, which we cannot separate with the present data. The quadrants are image quadrants and not anatomically registered ETDRS fields, so they may not capture the spatial pattern that the clinical rules concern; and the grading label of DDR may depend mainly on the amount and type of lesions rather than on their spread, for example because the boundaries between the middle grades are drawn on counts. The result also means that our rule $R_3$, which uses quadrant-wise haemorrhage evidence, is less well supported than the other rules.

\paragraph{Lesion classes contribute unequally.} Removing haemorrhages costs 0.030 QWK, removing microaneurysms 0.024, removing soft exudates 0.019, and removing hard exudates only 0.004; the last difference is negligible although exudates are the best-segmented class (Section~\ref{sec:res-seg}), presumably because their information is redundant with the other classes and with the image embedding (large exudates are also visible to the ImageNet features). When a single class is kept, microaneurysm evidence reaches 0.704, haemorrhage evidence 0.687, soft-exudate evidence 0.674 and exudate evidence 0.630; the last is no better than the embedding-only models (0.638--0.644; Table~\ref{tab:grading}). It is striking that microaneurysms, the class with the lowest segmentation quality (AUPR at most 0.17), are the most informative \emph{single} class. A plausible reason is that microaneurysms are the earliest lesion and that ICDR grade~1 is defined by their presence alone, so that even a noisy microaneurysm signal separates grade~0 from higher grades. This implies that better microaneurysm segmentation, which requires higher-resolution input than we could afford, has the greatest potential to improve the grader. As with all ablations in this paper, these are three-seed estimates on a single backbone and should be read as indications.

\begin{table}[!htbp]
\centering
\caption{Evidence ablation: the full \mname\ head trained with parts of the evidence tensor removed (zones or lesion classes set to zero everywhere, including in the rule scores). Three seeds (0--2), mean $\pm$ SD, leak-controlled test set. $\Delta$ is the mean QWK difference to the reference. With three seeds and SDs of 0.005--0.02, differences below about 0.01 are not distinguishable from noise.}
\label{tab:evablate}
\small
\resizebox{\textwidth}{!}{%
\begin{tabular}{@{}lcccc@{}}
\toprule
Evidence kept & QWK & $\Delta$QWK & Acc. & ECE \\
\midrule
Full evidence (reference) & 0.745$\pm$0.007 & -- & 0.748$\pm$0.009 & 0.048$\pm$0.009 \\
Whole-field zone only & 0.736$\pm$0.013 & -0.008 & 0.747$\pm$0.013 & 0.056$\pm$0.012 \\
Quadrants + centre (no whole field) & 0.733$\pm$0.022 & -0.012 & 0.742$\pm$0.010 & 0.055$\pm$0.024 \\
Without MA & 0.721$\pm$0.002 & -0.023 & 0.739$\pm$0.000 & 0.055$\pm$0.002 \\
Without HE & 0.715$\pm$0.011 & -0.029 & 0.729$\pm$0.006 & 0.058$\pm$0.010 \\
Without EX & 0.741$\pm$0.011 & -0.003 & 0.747$\pm$0.003 & 0.054$\pm$0.021 \\
Without SE & 0.726$\pm$0.009 & -0.019 & 0.736$\pm$0.008 & 0.032$\pm$0.001 \\
MA only & 0.704$\pm$0.018 & -0.041 & 0.717$\pm$0.006 & 0.062$\pm$0.007 \\
HE only & 0.687$\pm$0.004 & -0.058 & 0.716$\pm$0.005 & 0.067$\pm$0.017 \\
EX only & 0.630$\pm$0.003 & -0.115 & 0.680$\pm$0.003 & 0.083$\pm$0.009 \\
SE only & 0.674$\pm$0.009 & -0.071 & 0.697$\pm$0.001 & 0.076$\pm$0.009 \\
\bottomrule
\end{tabular}}
\end{table}

